\section{Adaptive standard-ray enumeration from minimum envelopes}
\label{sec:minimum-envelopes-native}

Report 74 replaces a particular over-refinement in the native sector
enumerator. The older standard phase formed every pairwise equality of
triangle potentials, lifted the resulting nonnegative directions, and checked
their full standard extremality. A crossing between two potentials above
the minimum can create a spurious direction. At quadrilateral matching
nullity at most two, the new method enumerates only changes of the actual
minimum and the feasible section endpoints. These are exactly the standard
rays, so their rank rejection stage disappears.

Three delivered production modules and their independent tests are now
integrated unchanged. The native adaptation additionally selects this method
automatically at nullity one or two and threads an explicit method option
through supplied-sector discovery. Higher nullity retains the complete
arrangement/support choice; the quadrilateral phase retains its old behavior.
This is a structural adaptive choice with a completeness proof, rather than
a heuristic that drops difficult candidates.

\subsection{Canonical lifting is linear on the right faces}

For a compatible allowed quadrilateral sector with $k\le t$ types, let
$C=\{q\ge0:Aq=0\}$ be its quadrilateral matching cone. The triangle graph
gives linear corner potentials $\ell_c(q)$, grouped by global vertex $v$.
Triangle coordinates have the form $\alpha_v+\ell_c(q)$. Their canonical
lift is
\[
 \Psi(q)_c=\ell_c(q)-m_v(q),\qquad
 m_v(q)=\min_{c\in G_v}\ell_c(q).
\]
Every standard sector vector is uniquely $\Psi(q)+\sum_v\beta_v L_v$,
$\beta_v\ge0$, with $L_v$ its vertex-link vector. A non-link extreme ray
must have all $\beta_v=0$, since a nonzero link summand has zero quadrilateral
projection and cannot be proportional to the whole vector.

Choose a minimum anchor $b_v$ in each group and define
\[
 F_b=\{q\in C:\ell_c(q)\ge\ell_{b_v}(q)\text{ for all }c\in G_v\}.
\]
On $F_b$ the lift is linear. Projection and lifting are inverse maps between
$F_b$ and the full standard coordinate face obtained by setting all selected
triangle anchors to zero. This is a face because every standard coordinate
is nonnegative. Extreme rays of a face are extreme rays of the whole cone.
Conversely every non-link standard extreme ray is canonical and belongs to
one such anchored face. Therefore its projected direction is a ray of the
\emph{minimum} subdivision, not an arbitrary pairwise-potential arrangement.

This face correspondence holds in every dimension. Within an anchor cell,
nonnegative summands of a lifted vector must keep every zero anchor zero;
global nonadditivity of $\Psi$ creates no extremality gap. In higher dimension
the minimum subdivision is the common lower normal fan of the potential
coefficient polytopes, equivalently that of their Minkowski sum, intersected
with $C$. This description does not bound arbitrary-dimensional fan size.
Canonical conversion itself is classical; the manuscript retains its
attribution to normal-surface conversion work.

\subsection{Complete nullity-two enumeration and its output bound}

Let $d=\dim_{\mathbb Q}\ker A\le2$. Normalize by $\sum q_i=1$, which meets
every nonzero nonnegative ray exactly once. The section is empty, a point,
or a bounded interval. For a nondegenerate interval choose exact rational
$u,w$ with $Au=Aw=0$, $\sum u_i=1$, $\sum w_i=0$, $w\ne0$, and write
$q(s)=u+sw$. Every coordinate inequality gives an exact lower/upper bound
or a constant feasibility test, producing $[L,U]$.

For each vertex group put $h_v(s)=\min_c\ell_c(u+sw)$. This is the lower
envelope of affine lines. The complete ray list consists of primitive
canonical lifts at $L,U$ and every genuine interior slope change of any
$h_v$. Between those parameters the same anchor realizes each minimum,
so there is a single linear minimum cell and no interior extreme direction.
At a genuine change the section is a vertex of an anchored face. An interior
zero of a nonnegative affine coordinate must be an identically zero
coordinate, so no additional quadrilateral event is missing. Parallel or
coincident lines and endpoint-only ties create no extra ray.

Dimension zero gives no non-link ray; dimension one is empty or a point.
A two-dimensional matching space may still have an empty or singleton
nonnegative section. All these cases are implemented explicitly. The raw
linear nullity is not confused with feasible cone dimension.

\begin{proposition}[Linear ray inventory]
For a nondegenerate section, let $p$ be retained triangle classes, $g$ their
vertex groups and $a\le4k$ the nonzero-labelled triangle-graph edges. Then
\[
 R\le p-g+2\le a-k+4\le3k+4.
\]
\end{proposition}
\begin{proof}
A line is minimal on an interval, since each comparison with another line
defines an interval. Concavity makes successive distinct envelope slopes
strictly decrease. A group with $p_v$ classes therefore has at most $p_v-1$
genuine changes; unioning changes and adding endpoints gives $p-g+2$.
After zero-label contraction, cycle consistency has rank at most
$a-p+g$. Its matching nullity two gives rank $k-2$, so
$k-2\le a-p+g$ and $p-g+2\le a-k+4$. Each allowed quad labels at most
four face-matching edges, hence $a\le4k$.
\end{proof}

The constant three is an upper bound, not a claimed sharp geometric example.
The local $p-g+2$ bound is often tighter. An abstract affine-potential family
can attain the corresponding envelope count, but that does not realize every
such potential system as a knot exterior.

The hull implementation sorts slopes, keeps the smallest intercept at equal
slope, and pops a previous line when the next intersection is not later than
its start. Exact fractions retain all genuine ties. Projected potential
intercepts/slopes are computed once and reused across all lifts. With
$p=O(k)$ and $R=O(k)$ the post-kernel arithmetic cost is
\[
 O(pk+p\log(p+1)+R(t+p+k))=O(tk),\qquad1\le k\le t.
\]
Kernel construction, source preparation and full coordinate validation are
charged separately. Coefficient sizes give polynomial bit time; this is not
a quadratic bit-time bound. The growing family has exponentially many
pieces but polynomially encoded output, so numeric piece counts cannot be
treated as input lengths.

\subsection{Independent certificates of complete coverage}

The new \code{normal-sector-envelope-v1} binds the exact triangulation,
allowed types, matching and section dimensions, rational section vectors,
endpoints, active-envelope corner records and the complete primitive ray list.
Rationals use reduced numerator/positive-denominator pairs. The coverage
checker imports neither the producer, its hull routine nor its contracted
kernel builder. It reconstructs an uncontracted corner graph from the actual
matching rows, derives its own potentials and independently checks nullity.

Supplied section vectors are checked against every matching equation and
normalization; the coordinate bounds certify the entire nonnegative section.
For an envelope, active slopes must strictly decrease and successive
intersections must occur in order within the section. Each active line comes
from an actual source corner. An omitted source line $f$ dominates the
proposed concave envelope $h$ exactly when the convex function $f-h$ is
nonnegative at its minimum. A supplied slope bracket identifies that minimum
at a breakpoint or endpoint, giving one exact dominance comparison per
corner rather than testing all corners at all candidate parameters.

These checks certify the entire minimum subdivision and then every lifted
ray, with no missing or duplicate output. A positive disc certificate is
still a separate native topology/essentiality proof. Coverage of this
supplied sector does not authenticate a diagram exterior or prove that no
other sector contains a useful disc.

Automatic standard enumeration now dispatches to envelopes for $d=1,2$;
empty $d=0$ auto results retain their previous empty schedule. Explicit
\code{method='envelope'} accepts only the standard phase with $d\le2$.
Explicit arrangement/supports remain available; higher-dimensional auto
selection still uses the old predicted-work comparison. The quadrilateral
phase does not silently become a standard phase.

The envelope's \code{max\_bases} counts begun actual output lifts, not
spurious linear systems. A partial iterator raises \code{SearchLimit} and
establishes no absence claim. The complete-coverage producer prechecks
\code{max\_rays} against the true planned output and returns inconclusive
without a partial certificate. Supplied-sector discovery exposes the method
selector and preserves independently checked positive/exhaustion protocols.
Its standard exhaustion says only that no vertex disc was found in that
sector. Caller cancellation remains cooperative throughout planning and
verification.

\subsection{All-size capped family and its component classification}

The supplied fixture attaches one tetrahedral ball to face two of a Fibonacci
layered solid torus, allowing the meridional quad in each original tetrahedron
and type one in the cap. Put $F_0=0$, $F_1=F_2=1$,
$a=F_{n+1}$, $b=F_{n+2}$, and for integers $x,y\ge0$ let
$h=\max(0,y-ax)$. Its canonical vector has original rows
\[
 (F_{n-i+1}x+h,F_{n-i+1}x+h,h,h;0,0,F_{n-i}x),\quad0\le i<n,
\]
and cap row
\[
 (ax+h-y,bx+h,0,h;0,y,0).
\]
The terminal face gluing forces the last meridional quad to be $x$,
the penultimate one to be $x$ when it exists, and the remaining quads
to satisfy $q_i=q_{i+1}+q_{i+2}$. Triangle matching makes the two lower
triangle entries a common value $\alpha$ and the two upper entries
$F_{n-i+1}x+\alpha$. The cap equations give triangle entries
$ax+\alpha-y,bx+\alpha,\beta_1,\alpha$. Nonnegativity is exactly
$x,y\ge0$, $\alpha\ge h$ and $\beta_1\ge0$. Subtracting the old
vertex-link minimum $\beta_0=\alpha-h$ gives the displayed canonical form. There
are only two minimum cells, separated by $y=ax$, so the non-link standard
ray parameters are $(1,0),(1,a),(0,1)$. The middle ray is absent from the
two quadrilateral endpoints, although one endpoint also supplies an
essential disc. This demonstrates stronger complete enumeration, not a
claim that Q screening fails all useful recognition tests.

The README asserts a component formula; its delivered family script checks
finitely many parameters. The following geometric argument records the
uniform classification explicitly instead of treating those tests as proof.

\begin{proposition}[Every integral vector in the capped sector]
The canonical vector above consists of $x+h$ discs, exactly $x$ of which
have essential boundary. Adding nonnegative vertex-link multiplicities
$\beta_0,\beta_1$ gives $x+h+\beta_0+\beta_1$ disc components and the same
essential count $x$. These parameters describe all integral vectors of this
fixed sector for every $n\ge1$.
\end{proposition}
\begin{proof}
Restriction to the original torus is $xD_n+hL$, where $D_n$ is the primitive
meridian of Section~\ref{sec:unit-ray-native} and $L$ its boundary vertex
link. It is a disjoint union of $x$ meridian discs and $h$ inessential discs.
Each normal piece in a tetrahedron meets a given face in at most one arc.
The only type missing that face is its opposite triangle, here cap type two,
whose coordinate is zero. Thus every cap piece attaches along exactly one
disjoint boundary interval to one old component. The incidence graph for
each old component is a star with disc leaves: gluing preserves its disc
type and cannot merge two old components or create a separate one.

The cap is the fixture's boundary retriangulation. Its new outer three-face
disk maps to the old attaching triangle by a boundary-fixing subdivision
homeomorphism. Each added disc's outer boundary arc maps to a path homotopic
relative endpoints to its attaching arc. Hence the resulting boundary circles
retain their old peripheral classes, preserving the $x$ essential circles.
This argument takes place away from identified boundary vertices and respects
the formal face gluings.

The two source vertices are boundary vertices. Their added vertex links are
inessential disc components. Canonical decomposition separates these links
from every standard vector. The terminal meridional quad has coefficient
$F_1x=x$ and the cap quad has coefficient $y$, so integrality forces integer
$x,y$; their canonical triangle values and link minima are then integral.
Matching and nonnegativity give exactly the displayed parameters and their
nonnegative link additions.
\end{proof}

Consequently primitivity alone does not imply connectedness: $x=2,y=1$
has coordinate gcd one but two essential discs. The supported-ray theorem
requires its additional nullity-one hypothesis. On this family the previous
arrangement method admits $n+2$ positive directions and rejects $n-1$;
the envelope method emits exactly three. Controls with cap type zero have
two rays and are retained separately.

\subsection{Native validation and isolated timings}

All 1,351 maintained test methods pass in 224.903 seconds. The focused
42-method run passes in 3.481 seconds, including the unchanged native
unit-ray disc tests. New integration tests forbid calls to the pairwise
hyperplane generator and standard-rank filter during automatic envelope
enumeration, and independently replay positive and negative sector results.
They also exercise actual-ray resource caps, the extra non-Q disc, strict
method/phase validation and rejected missing coverage rays.

The native audit selects 9,995 sectors across all 48 frozen triangulations
and audits all 9,344 with matching nullity at most two. The declared
selection includes every support on the small triangulations and prescribed
ray/small/random supports on the larger ones; it is not every possible
sector on every larger triangulation. Every emitted set matches the frozen
complete standard-ray list filtered by its allowed types. There are 4,853
ray occurrences and 804 interior-ray occurrences across those repeated
queries. Automatic and explicit envelope output agree in 9,360 calls,
including fifteen capped-family configurations and one extra-disc example.
Fifteen fresh Regina full enumerations also match. All 567 captured source
hashes remain unchanged during the audit and timings. Fresh family controls
check nine ray lists through base size 128 and 97 integral surface vectors
against the maintained component observer, supporting the uniform cap proof.

The delivered complete-coverage evidence separately replays sixty eligible
independent-audit certificates, six coverage-benchmark certificates and one
essential-disc certificate with enumeration imports blocked. Six
higher-nullity model-only records remain outside this API. Native publication
checks replay those saved coverage proofs with the current independent
checker, together with the new capacity certificates. Empty and singleton
sections, missing rays, false breakpoints, slope brackets, source changes,
noncanonical fractions and callback cancellation have explicit controls.

The isolated benchmark uses five interleaved four-arm rounds and one warmup
per arm: 260 measured calls and 52 warmups, all complete. Its reference is
the frozen \code{654c0b2a7} native sector code. Both arms pay full kernel
construction. Enumeration times include all emitted coordinates, their
serialization and digest; this existing iterator API has no portable
coverage proof, so independent coverage replay is not claimed inside that
timing. The separate discovery times include its native essential-disc or
exhaustion proof, independent replay and full proof serialization. Fixtures
are constructed before timing, and ratios are paired medians.

\begin{samepage}
\noindent\textbf{Complete standard-ray enumeration and output.}
\begin{center}
\small
\begin{tabular}{lrrrrr}
\toprule
Input & old ms & new ms & old/new & old A/A & new A/A\\
\midrule
Base 1, cap 1 & 1.068 & 0.714 & 1.496 & 0.959 & 0.985\\
Base 4, cap 1 & 12.053 & 1.625 & 7.498 & 0.986 & 1.018\\
Base 8, cap 1 & 79.621 & 3.624 & 21.642 & 0.979 & 1.005\\
Base 16, cap 1 & 741.848 & 10.808 & 68.440 & 1.016 & 1.013\\
Base 32, cap 1 & 8449.232 & 43.139 & 197.217 & 1.005 & 0.996\\
Base 8, cap 0 & 26.903 & 3.304 & 8.209 & 1.002 & 0.995\\
Base 16, cap 0 & 137.210 & 9.984 & 13.791 & 0.986 & 0.997\\
\bottomrule
\end{tabular}
\end{center}

\end{samepage}

At base size 32, full enumeration drops from about 8.449 seconds to 0.043
seconds, a paired 197.22-fold gain, returning the same three rays. The old
method generates 34 positive directions and rejects 31. Gains are smaller
on the smallest source and substantial on the cap-zero controls without a
genuine minimum break, where projected hyperplane/rank work is still avoided.
These are local enumeration measurements on genuine finite triangulations.

\begin{samepage}
\noindent\textbf{Complete supplied-sector disc discovery and replay.}
\begin{center}
\small
\begin{tabular}{lrrrrr}
\toprule
Input & old ms & new ms & old/new & old A/A & new A/A\\
\midrule
Base 1, cap 1 & 1.743 & 1.484 & 1.166 & 0.985 & 1.015\\
Base 4, cap 1 & 7.796 & 3.375 & 2.301 & 1.006 & 1.000\\
Base 8, cap 1 & 28.471 & 6.557 & 4.339 & 0.999 & 1.006\\
Base 16, cap 1 & 144.934 & 16.111 & 9.002 & 0.993 & 0.986\\
Base 32, cap 1 & 918.508 & 53.136 & 17.283 & 1.005 & 1.002\\
Empty sector & 0.626 & 0.627 & 0.996 & 0.997 & 0.999\\
\bottomrule
\end{tabular}
\end{center}

\end{samepage}

Complete discovery at base size 32 drops from about 0.919 to 0.053 seconds,
a paired 17.28-fold gain. The empty-sector control is near parity. These
smaller ratios include the remaining kernel, geometry and independent
positive-proof costs and are not substituted by the larger enumeration
ratios. All complete discovery arms return the same status; individual
witnesses are checked, rather than assumed equal from status alone.
A/A medians range from 0.959 to 1.018.

Four separately timed new-only calls produce and replay complete coverage
certificates at base sizes 1, 8, 32 and 64. Their full elapsed times are
approximately 0.002, 0.010, 0.135 and 0.765 seconds, with certificate sizes
501, 1,067, 3,915 and 9,985 bytes. No speedup is inferred from these
unpaired capacity observations. Source/coverage replay and asymptotic
proofs remain separate evidence from wall-clock measurements.

\begin{sloppypar}
Reproduction is the maintained \code{normal\_orbit\_research.envelopes}
driver, with \code{audit --fresh-regina} or \code{benchmark --rounds 5}
and an explicit \code{--output} path, from \code{fast/}. The raw audit,
timings, test logs, source hashes and publication checks are retained in
\code{data/envelopes-*}. The complete original report and all its proof
appendices remain in \code{reports/74}.
\end{sloppypar}


The global obligation remains selection of a complete, sufficiently small
family of useful sectors, with their source and transition costs. Matching
nullity of an allowed quad sector is different from full-coordinate nullity
on the positive support of a discovered standard ray. The latter being one
does not prove that the former is at most two. Report 74 also proves that
checked Farkas zero-coordinate implications can restrict a cone to an
equivalent smaller face, but supplies no universal reduction to this
low-dimensional class. This native local improvement and the preserved
cycle/region theorems therefore advance specific operations while leaving
general QP recognition unproved.
